\documentclass[11pt, a4paper]{article}

\usepackage[T1]{fontenc}
\usepackage[utf8]{inputenc}
\usepackage{tgpagella}
\usepackage[spanish, es-noshorthands]{babel}
\usepackage{amsmath, amssymb, amsfonts}
\usepackage{graphicx}
\usepackage{booktabs}
\usepackage{tabularx}
\usepackage[table, xcdraw]{xcolor}
\usepackage[most]{tcolorbox}
\usepackage[american]{circuitikz}
\usepackage[margin=2.5cm]{geometry}
\usepackage{fancyhdr}

\pagestyle{fancy}
\fancyhf{}
\setlength{\headheight}{16pt}
\lhead{\textbf{UCB Sede Tarija} | Ing. Mecatrónica}
\chead{Introducción a la Mecatrónica}
\rhead{Semana 5 - Sesión 2}
\lfoot{Prof: M.Sc. Ing. Bernardo Quiroga Turdera}
\rfoot{Página \thepage}
\renewcommand{\headrulewidth}{0.4pt}

\definecolor{ucbblue}{RGB}{0, 51, 102}

\begin{document}

\begin{tcolorbox}[colback=ucbblue!5!white, colframe=ucbblue, title=\textbf{Hoja de Referencia y Derivaciones: Amplificadores Operacionales}]
    \centering \large \textbf{Unidad 2 - Configuraciones Clásicas}
\end{tcolorbox}

\section*{1. Modelo Ideal del Amplificador Operacional}
El Op-Amp ideal simplifica el análisis de circuitos asumiendo las siguientes propiedades cuando opera en su región lineal y bajo realimentación negativa:
\begin{itemize}
    \item Ganancia en lazo abierto infinita: $A \rightarrow \infty$.
    \item Impedancia de entrada infinita: $R_{in} \rightarrow \infty$.
    \item Impedancia de salida nula: $R_{out} \rightarrow 0$.
\end{itemize}

\textbf{Consecuencias operativas para la resolución de circuitos:}
\begin{align}
    i_+ &= i_- = 0 \quad \text{(Corrientes de entrada nulas)} \\
    V_+ &\approx V_- \quad \text{(Cortocircuito virtual debido a la realimentación negativa)}
\end{align}

\section*{2. Amplificador Inversor}
\begin{minipage}{0.45\textwidth}
    \begin{circuitikz}[scale=0.8, transform shape, thick]
        \draw (0,0) node[op amp] (opamp) {};
        \draw (opamp.+) -- ++(0,-0.5) node[ground]{};
        \draw (opamp.-) node[above]{$V_-$} to[R, l_=$R_{in}$] ++(-2.5,0) node[left]{$V_i$};
        \draw (opamp.-) -- ++(0,1.5) coordinate(tmp) to[R, l=$R_f$] (tmp -| opamp.out) -- (opamp.out) to[short, -o] ++(0.5,0) node[right]{$V_o$};
    \end{circuitikz}
\end{minipage}
\hfill
\begin{minipage}{0.5\textwidth}
    \textbf{Esqueleto de Derivación:}
    \begin{align*}
        V_+ &= 0 \implies V_- = 0 \\
        i_{in} &= i_f \implies \frac{V_i - V_-}{R_{in}} = \frac{V_- - V_o}{R_f}
    \end{align*}
    Sustituyendo la tierra virtual ($V_-=0$):
    \begin{equation*}
        \boxed{ A_v = \frac{V_o}{V_i} = -\frac{R_f}{R_{in}} }
    \end{equation*}
\end{minipage}

\section*{3. Amplificador No Inversor}
\begin{minipage}{0.45\textwidth}
    \begin{circuitikz}[scale=0.8, transform shape, thick]
        % Se utiliza "noinv input up" para colocar V+ arriba y V- abajo, 
        % rutando la realimentación por debajo y evitando cruces de líneas.
        \draw (0,0) node[op amp, noinv input up] (opamp) {};
        \draw (opamp.+) to[short, -o] ++(-1.5,0) node[left]{$V_i$};
        \draw (opamp.-) node[above]{$V_-$} to[short] ++(0,-1.5) coordinate(tmp) to[R, l_=$R_f$] (tmp -| opamp.out) -- (opamp.out) to[short, -o] ++(0.5,0) node[right]{$V_o$};
        \draw (opamp.-) -- ++(-0.5,0) to[R, l_=$R_g$] ++(0,-1.5) node[ground]{};
    \end{circuitikz}
\end{minipage}
\hfill
\begin{minipage}{0.5\textwidth}
    \textbf{Esqueleto de Derivación:}
    \begin{align*}
        V_+ &= V_i \implies V_- = V_i \\
        V_- &= V_o \left( \frac{R_g}{R_f + R_g} \right) \text{ (Divisor)}
    \end{align*}
    Sustituyendo e igualando a $V_i$:
    \begin{equation*}
        \boxed{ A_v = \frac{V_o}{V_i} = 1 + \frac{R_f}{R_g} }
    \end{equation*}
\end{minipage}

\section*{4. Tabla Comparativa de Topologías}
\begin{table}[htbp]
    \centering
    \renewcommand{\arraystretch}{1.5}
    \begin{tabularx}{\textwidth}{@{} l X X @{}}
    \toprule
    \rowcolor[HTML]{EFEFEF}
    \textbf{Parámetro Analítico} & \textbf{Amplificador Inversor} & \textbf{Amplificador No Inversor} \\ \midrule
    Relación de Fase (Signo) & Invertida ($-$) & En fase ($+$) \\
    Expresión de Ganancia $A_v$ & $-R_f / R_{in}$ & $1 + R_f / R_g$ \\
    Impedancia de Entrada ($Z_{in}$) & $\approx R_{in}$ (Carga la fuente) & Idealmente $\infty$ (No carga) \\
    Aplicación Principal & Sumadores, inversores de señal & Buffers, adaptación de impedancias \\
    \bottomrule
    \end{tabularx}
\end{table}

\end{document}
